% Article: Polygamma values at complex multiplication points via Eisenstein row-sums
% - Created (UTC): 2026-06-11T13:07:42Z
% - Repository HEAD: 68236dab4fa7979b44b115e72bd7c124ae550d92
% Source corpus: src/PolyLog/docs/reports/polygamma-complex-arguments-cm-lattices__d41f7be20c61.md
% and the commit range 272af472..68236dab4 of VladimirReshetnikov/Smithereens.
\documentclass[11pt]{article}
\usepackage[margin=1.1in]{geometry}
\usepackage{amsmath,amssymb,amsthm}
\usepackage{booktabs}
\usepackage{microtype}
\usepackage{xcolor}
\usepackage[colorlinks=true,linkcolor=blue!50!black,citecolor=blue!50!black,urlcolor=blue!60!black]{hyperref}
\usepackage{xurl}

\newtheorem{result}{Result}[section]
\newtheorem{law}[result]{Law}
\newtheorem{lemma}[result]{Lemma}
\newtheorem{negative}[result]{Negative result}
\theoremstyle{remark}
\newtheorem{remark}[result]{Remark}

\newcommand{\Li}{\operatorname{Li}}
\newcommand{\Cl}{\operatorname{Cl}}
\newcommand{\sech}{\operatorname{sech}}
\newcommand{\dps}[1]{\textsf{\footnotesize[#1]}}

\title{Polygamma Values at Complex Multiplication Points:\\
Eisenstein Row-Sums, Chowla--Selberg Periods,\\ and Hilbert Class Polynomials}
\author{Vladimir Reshetnikov\thanks{%
v.reshetnikov@gmail.com. Prepared with agentic assistance (Anthropic Claude) in the
\emph{Smithereens} repository,
\texttt{github.com/VladimirReshetnikov/Smithereens}, \texttt{src/PolyLog}.
Provenance: created 2026-06-11 (UTC); repository snapshot
\texttt{68236dab4fa7979b44b115e72bd7c124ae550d92}.}}
\date{June 11, 2026}

\begin{document}
\maketitle

\begin{abstract}
Summing an Eisenstein series of the lattice $\mathbb{Z}+\mathbb{Z}\tau$ row by row turns it
into a series of polygamma values at the non-real points $n\tau$. At complex multiplication
(CM) points the Eisenstein values are pinned by classical theory, and the row-sum identity
converts that information into explicit evaluations of infinite sums of polygamma values at
complex arguments. We record three layers of consequences, all verified numerically at
$50$--$300$ digits with PSLQ-canary discipline. First, \emph{vanishing identities}: at the
square and hexagonal lattices the missing modular forms collapse polygamma sums to rational
multiples of $\zeta(4)$ and $\zeta(6)$. Second, \emph{period identities}: the surviving
Eisenstein values make three distinct Chowla--Selberg periods --- $\Gamma(1/4)$, $\Gamma(1/3)$,
and $\Gamma(1/8)\Gamma(3/8)$ --- visible as polygamma sums, and the full Eisenstein tower over
both maximal-symmetry lattices is polygamma-explicit through weight~$12$. Third, \emph{class
field theory}: for discriminants of class number $2$ through $5$ ($-20$, $-15$, $-23$, $-39$,
$-47$), cross-class ratios are algebraic numbers whose norms forecast the prime content of
class-product period monomials, and the row-sum values recover the Hilbert class polynomials
$H_{-39}$ and $H_{-47}$ outright --- integer coefficients of $20$ and more digits emerging from
sums of polygamma values at quadratic-irrational complex points. A separate elementary layer
(reflection composed with conjugation) explains MathWorld's ``surprising'' FoxTrot digamma
identity as one point of a two-parameter family. We also record the honest negatives: the
components opaque to these mechanisms are Eichler-integral/quasi-modular data, and
cross-class \emph{sums} of weight-$4$ values are not Galois-clean.
\end{abstract}

\section{Introduction}

The starting question was deliberately open-ended: \emph{are there interesting identities for
polygamma values at non-real arguments --- or for a single real or imaginary component of such
values, even when the other component is opaque?}

The mechanism that answers it is elementary. For $a\in\mathbb{C}$ off the real axis and
$s\ge 1$, splitting the bilateral sum at $m=0$ and using
$\psi^{(s)}(a)=(-1)^{s+1}s!\sum_{m\ge 0}(m+a)^{-s-1}$ gives
\begin{equation}\label{eq:bilateral}
\sum_{m\in\mathbb{Z}}\frac{1}{(m+a)^{s+1}}
=\frac{(-1)^{s+1}\,\psi^{(s)}(a)+\psi^{(s)}(1-a)}{s!}
\;=\;\frac{\psi^{(s)}(a)+\psi^{(s)}(1-a)}{s!}\quad(s\text{ odd}).
\end{equation}
Thus \emph{row-summing a lattice sum turns it into a polygamma sum}. For the weight-$2k$
Eisenstein series of $\mathbb{Z}+\mathbb{Z}\tau$ ($\operatorname{Im}\tau>0$, $k\ge 2$),
summing first over $m$ in each row $n$ and noting that the rows $\pm n$ coincide (the power is
even),
\begin{equation}\label{eq:rowsum}
G_{2k}(\tau)\;=\;\sideset{}{'}\sum_{(m,n)\in\mathbb{Z}^2}\frac{1}{(m+n\tau)^{2k}}
\;=\;2\zeta(2k)+\frac{2}{(2k-1)!}\sum_{n\ge 1}
\Bigl[\psi^{(2k-1)}(n\tau)+\psi^{(2k-1)}(1-n\tau)\Bigr].
\end{equation}
Everything in this article is an exploitation of \eqref{eq:rowsum} at points $\tau$ where the
left-hand side is known: the maximal-symmetry CM points $\tau=i$ and $\tau=\rho=e^{i\pi/3}$,
the isogenous neighbour $2i$ and the discriminant~$-8$ point $i\sqrt2$, and the
reduced-form points of the imaginary-quadratic discriminants of class numbers $2$--$5$.

\paragraph{Epistemic status.}
These are results of experimental mathematics, recorded in the PolyLog corpus of the
Smithereens repository. Unless explicitly marked classical, each identity was discovered or
confirmed by PSLQ over a deliberately prepared basis and verified with \texttt{mpmath} at the
stated precision (residuals quoted per identity, typically $50$--$300$ significant digits),
with the project's \emph{canary} discipline guarding against PSLQ false positives: every
candidate relation must continue to hold when re-evaluated at a substantially higher precision
than the search precision. Polygamma values at complex $a$ were computed via the
Hurwitz-zeta shift $\psi^{(m)}(a)=(-1)^{m+1}m!\,(a^{-m-1}+\zeta(m+1,1+a))$, which keeps the
Hurwitz zeta in its numerically safe region. None of the new identities below currently
carries a symbolic proof certificate; converting them into the repository's certificate
format is future work (the certificate program is the subject of a companion article).

\section{The elementary layer: reflection composed with conjugation}\label{sec:elementary}

Before any CM input, one family of components is fixed by symmetry alone. The Schwarz
reflection $\psi^{(m)}(\bar z)=\overline{\psi^{(m)}(z)}$ turns the reflection law into a
statement about a single real component, leaving the other component untouched.

\begin{result}[vertical-line trigamma law]\label{res:line}
On the line $\operatorname{Re}z=\tfrac12$ the reflection law alone evaluates the real part of
the trigamma function at every height: for all real $t$,
\[
\operatorname{Re}\psi^{(1)}\!\Bigl(\tfrac12+it\Bigr)=\frac{\pi^2}{2\cosh^2(\pi t)} .
\]
In particular $\operatorname{Re}\psi^{(1)}\!\bigl(\tfrac{1+i}{2}\bigr)=\pi^2/(2\cosh^2(\pi/2))$
\dps{residual $0$ at $50$ dps}. Off the symmetric line, combining reflection, conjugation and
the recurrence gives e.g.
\[
\operatorname{Re}\psi^{(1)}(i)=-\frac12-\frac{\pi^2}{2\sinh^2\pi}
\qquad\text{\dps{residual $0$ at $50$ dps}.}
\]
The proof of the first identity is one line: on $\operatorname{Re}z=\tfrac12$ one has
$1-z=\bar z$, so the reflection law
$\psi^{(1)}(z)+\psi^{(1)}(1-z)=\pi^2/\sin^2(\pi z)$ reads
$2\operatorname{Re}\psi^{(1)}(z)=\pi^2/\sin^2(\pi z)$ there.
\end{result}

The imaginary parts on this line stay opaque --- they are odd-row lattice data in the sense of
Section~\ref{sec:opaque}.

\subsection{The FoxTrot identity demystified}

MathWorld's \emph{Digamma Function} page \cite{MathWorldDigamma} calls the following
identity ``surprising''; it arises from the FoxTrot series $\sum(-1)^{n-1}n^2/(n^3+1)$,
whose partial fractions hit the roots of $n^2-n+1$:
\begin{equation}\label{eq:foxtrot}
-\psi\bigl(\tfrac{\omega}{2}\bigr)-\psi\bigl(-\tfrac{\omega^2}{2}\bigr)
+\psi\bigl(\tfrac{1+\omega}{2}\bigr)+\psi\bigl(\tfrac{1-\omega^2}{2}\bigr)
=2\pi\,\sech\Bigl(\tfrac{\sqrt3\,\pi}{2}\Bigr),\qquad \omega=e^{i\pi/3}.
\end{equation}

\begin{law}[line law for digamma differences]\label{law:lineline}
For every shift $c\in(0,1)$, on the line $\operatorname{Re}z=(1-c)/2$,
\[
\operatorname{Re}\bigl[\psi(z+c)-\psi(z)\bigr]=\pi\,\operatorname{Re}\cot(\pi\bar z).
\]
\dps{verified for $c=\tfrac12,\tfrac13,\tfrac23$ at generic heights $t$}
\end{law}

The mechanism is the same one-liner: on that line the shift satisfies $z+c=1-\bar z$, so
reflection composed with conjugation eliminates the digamma values, leaving an elementary
$\sinh/\cos$ closed form. The FoxTrot identity \eqref{eq:foxtrot} is the instance
$c=\tfrac12$, $z=\tfrac14+i\tfrac{\sqrt3}{4}$: its four arguments pair into
$2\operatorname{Re}[\psi(z+\tfrac12)-\psi(z)]$, and on $\operatorname{Re}z=\tfrac14$ one has
$\pi\operatorname{Re}\cot(\pi\bar z)=\pi\sech(2\pi t)$. The identity is thus the point
$t=\sqrt3/4$ on a line of identities, itself one member of a two-parameter family --- no
complex multiplication required.

Differentiating in $t$ ladders the law up the polygamma tower
\dps{weights $2,3$ verified at $10^{-50}$}:
\begin{align*}
\operatorname{Im}\bigl[\psi^{(1)}(\tfrac34+it)-\psi^{(1)}(\tfrac14+it)\bigr]
&=2\pi^2\sech(2\pi t)\tanh(2\pi t),\\
\operatorname{Re}\bigl[\psi^{(2)}(\tfrac34+it)-\psi^{(2)}(\tfrac14+it)\bigr]
&=8\pi^3\sech^3(2\pi t)-4\pi^3\sech(2\pi t).
\end{align*}
At each weight the alternating component ($\operatorname{Re}$ at odd weight,
$\operatorname{Im}$ at even weight) stays opaque, exactly as in
Section~\ref{sec:opaque}.

\section{Maximal symmetry: vanishing identities and \texorpdfstring{$\Gamma$}{Gamma}-periods}
\label{sec:maximal}

At the two CM points with extra symmetry --- $\tau=i$ (square lattice, $j=1728$) and
$\tau=\rho=e^{i\pi/3}$ (hexagonal lattice, $j=0$) --- classical theory pins every $G_{2k}$
exactly. The square lattice admits no weight-$6$ form ($i\cdot(\mathbb{Z}+\mathbb{Z}i)
=\mathbb{Z}+\mathbb{Z}i$ forces $G_6(i)=i^{-6}G_6(i)=-G_6(i)$), and the hexagonal lattice no
weight-$4$ form; $G_4(i)$ and $G_6(\rho)$ are the lemniscatic and equianharmonic period
powers --- monomials in $\Gamma(1/4)$, resp.\ $\Gamma(1/3)$, by Chowla--Selberg \cite{ChowlaSelberg}.

\subsection{Vanishing identities}

\begin{result}[the CM gems]\label{res:vanishing}
\[
\sum_{n\ge 1}\Bigl[\psi^{(5)}(ni)+\psi^{(5)}(1-ni)\Bigr]=-120\,\zeta(6)
\qquad\text{\dps{residual $1.7\times10^{-49}$}},
\]
\[
\sum_{n\ge 1}\Bigl[\psi^{(3)}(n\rho)+\psi^{(3)}(1-n\rho)\Bigr]=-6\,\zeta(4),
\qquad \rho=e^{i\pi/3}
\qquad\text{\dps{residual $6.8\times10^{-50}$}}.
\]
\end{result}

Each individual bracket is a transcendental-looking complex number (elementary in
$\coth/\operatorname{csch}$ of $\pi n\sqrt3/2$ and the like after the bilateral cotangent
expansion, but nobody would guess the totals); the sums collapse to rational multiples of
$\zeta(6)$, $\zeta(4)$ \emph{because the corresponding modular forms do not exist}. Via
\eqref{eq:rowsum} these are exactly the statements $G_6(i)=0$ and $G_4(\rho)=0$.

\subsection{Chowla--Selberg made polygamma-visible}

\begin{result}[period identities]\label{res:periods}
\[
2\zeta(4)+\frac13\sum_{n\ge 1}\Bigl[\psi^{(3)}(ni)+\psi^{(3)}(1-ni)\Bigr]
=\frac{\Gamma(1/4)^8}{960\,\pi^2}
\qquad\text{\dps{residual $2.1\times10^{-50}$}},
\]
\[
2\zeta(6)+\frac1{60}\sum_{n\ge 1}\Bigl[\psi^{(5)}(n\rho)+\psi^{(5)}(1-n\rho)\Bigr]
=\frac{\Gamma(1/3)^{18}}{8960\,\pi^6}
\qquad\text{\dps{residual $2.0\times10^{-60}$}}.
\]
\end{result}

Polygamma values at purely imaginary (resp.\ hexagonal) points, summed over the integer
multiples of the point, know the lemniscate constant and its equianharmonic sibling. The
atoms $\Gamma(1/4)^2/\pi$ and $\Gamma(1/3)$ are precisely the CM-period currency of the
repository's hypergeometric reducer (\texttt{cmPeriodReduce} in \texttt{src/Hypergeometric}),
so these identities tie the polygamma world directly to that corpus.

\subsection{The complete low-weight table}

The symmetry argument and the classical Hurwitz recursion \cite{Hurwitz}
\[
(2n+1)(2n-1)(n-3)\,G_{2n}=3\sum_{m=2}^{n-2}(2m-1)(2n-2m-1)\,G_{2m}G_{2n-2m}
\]
extend both subsections through weight~$12$ with no new constants
\dps{all entries verified, residuals $10^{-60}$--$10^{-263}$}:

\begin{center}
\begin{tabular}{c l l}
\toprule
weight $2k$ & square lattice $\tau=i$ & hexagonal lattice $\tau=\rho$\\
\midrule
$4$ & $\Gamma(1/4)^8/(960\pi^2)$ & $\mathbf{0}$\\
$6$ & $\mathbf{0}$ & $\Gamma(1/3)^{18}/(8960\pi^6)$\\
$8$ & $\tfrac37 G_4(i)^2=3\,\Gamma(1/4)^{16}/(7\cdot 960^2\pi^4)$ & $\mathbf{0}$\\
$10$ & $\mathbf{0}$ & $\mathbf{0}$\\
$12$ & $\tfrac{1}{143}\bigl(18\,G_4(i)^3+25\,G_6(i)^2\bigr)$ type
     & $\tfrac{25}{143}G_6(\rho)^2\propto\Gamma(1/3)^{36}$\\
\bottomrule
\end{tabular}
\end{center}

Every bold zero is a vanishing identity
$\sum_{n\ge1}[\psi^{(w-1)}(n\tau)+\psi^{(w-1)}(1-n\tau)]=-(w-1)!\,\zeta(w)$, and every
non-zero entry is a $\psi^{(w-1)}$-sum evaluation. The pattern is exhaustive: $G_{2k}(i)=0$
iff $2k\not\equiv 0\pmod 4$, $G_{2k}(\rho)=0$ iff $2k\not\equiv 0\pmod 6$, and all non-zero
values are $\mathbb{Q}$-monomials in the two period powers --- so the \emph{entire} Eisenstein
tower over both lattices is polygamma-explicit.

\section{Beyond maximal symmetry: \texorpdfstring{$2i$ and $i\sqrt2$}{2i and i sqrt(2)}}
\label{sec:nextpoints}

The same row-sum at the next two CM points, with closed forms found by PSLQ
\dps{$120$-dps canaries}:

\begin{result}\label{res:2i}
\begin{gather*}
G_4(2i)=\frac{11}{15360}\,\frac{\Gamma(1/4)^8}{\pi^2}
\quad\text{\dps{canary $0$ at $120$ dps}},\\
G_4(i\sqrt2)=\frac{\Gamma(1/8)^4\,\Gamma(3/8)^4}{4608\,\pi^2}
\quad\text{\dps{canary $3.9\times10^{-121}$}}.
\end{gather*}
The weight-6 row at the same points \dps{$160$-dps canaries}:
\[
G_6(2i)=\frac{\Gamma(1/4)^{12}}{2^{14}\cdot 5\cdot\pi^3},\qquad
G_6(i\sqrt2)=\frac{\Gamma(1/8)^6\,\Gamma(3/8)^6}{2^{12}\cdot 3^3\cdot 5\cdot\pi^3}.
\]
\end{result}

The first value confirms the isogeny picture: $2i$ sits over the same field $\mathbb{Q}(i)$,
so only the rational factor changes ($G_4(i)/G_4(2i)=16/11$). The second imports a genuinely
\emph{new period}: discriminant $-8$, field $\mathbb{Q}(\sqrt{-2})$ --- the polygamma sum
$\sum_n[\psi^{(3)}(ni\sqrt2)+\psi^{(3)}(1-ni\sqrt2)]$ evaluates in $\Gamma(1/8)\Gamma(3/8)$.
The fundamental unit $1+\sqrt2$ was offered in the PSLQ basis and declined (exponent $0$).
Together with Section~\ref{sec:maximal} this makes three distinct Chowla--Selberg periods
polygamma-visible.

\section{Class number 2: discriminants \texorpdfstring{$-20$ and $-15$}{-20 and -15}}
\label{sec:h2}

At $\mathbb{Q}(\sqrt{-5})$ (class number $2$) a single-value logarithmic PSLQ \emph{must}
fail: the algebraic factor in $G_4/\Omega^4$ lives in the Hilbert class field and is not a
monomial. The structure that works is the Hecke/Damerell pair \cite{Damerell} over the two reduced forms
$x^2+5y^2$ ($\tau_1=i\sqrt5$) and $2x^2+2xy+3y^2$ ($\tau_2=(-1+i\sqrt5)/2$).

\begin{result}[disc $-20$]\label{res:d20}
\[
\frac{G_4(\tau_1)}{G_4(\tau_2)}=\frac{29+12\sqrt5}{44}=\frac{(3+2\sqrt5)^2}{44}
\quad\text{\dps{canary $0$ at $160$ dps}},
\]
\[
G_4(\tau_1)\,G_4(\tau_2)
=\frac{11\,\varphi^4\,\Gamma(3/20)^8\,\Gamma(7/20)^8}{2^{14}\cdot3^4\cdot5^4\cdot\pi^4}
\quad\text{\dps{canary $5.3\times10^{-160}$}},
\]
where $\varphi$ is the golden ratio. The ratio is algebraic of degree $2$ with norm
$N(29+12\sqrt5)=121=11^2$; the class product is a period monomial carrying the prime $11$
and a golden unit.
\end{result}

\begin{result}[disc $-15$]\label{res:d15}
With $\tau_1=(1+i\sqrt{15})/2$ for $x^2+xy+4y^2$ and the conjugate pair
$\tau_{2,3}=(\pm1+i\sqrt{15})/4$ representing $2x^2+xy+2y^2$, write
$\widehat G_4=\tfrac12\bigl[G_4(\tau_2)+G_4(\tau_3)\bigr]
=\operatorname{Re}G_4(\tau_2)$ for the class-pair average. Then
\dps{$160$-dps canaries}:
\[
\frac{G_4(\tau_1)}{\widehat G_4}=\frac{2\,(4+\sqrt5)^2}{77},\qquad
G_4(\tau_1)\,\widehat G_4
=\frac{77\,\bigl(\Gamma(\tfrac1{15})\Gamma(\tfrac2{15})\Gamma(\tfrac4{15})
\Gamma(\tfrac8{15})\bigr)^4}{2^{13}\cdot3^6\cdot5^5\cdot\pi^4}.
\]
Unlike disc $-20$, the non-principal value here is genuinely complex (the form
$(2,1,2)$ has $|\tau_2|=1$, where conjugation acts nontrivially on $G_4$); the
class-pair average is the Galois-natural object.
\end{result}

The $\Gamma$-block in Result~\ref{res:d15} is exactly the set of residues
$\{1,2,4,8\}$ with $\chi_{-15}=+1$; the golden unit is declined by PSLQ (exponent $0$); and
the prime content $77=7\cdot11$ of the product matches the ratio's denominator
($N(4+\sqrt5)=11$). The heuristic that emerges, used at every later discriminant:
\emph{run the ratio first; its norm names the primes the class-product basis needs.}

\section{Class number 3: discriminant \texorpdfstring{$-23$}{-23}}\label{sec:h3}

$\mathbb{Q}(\sqrt{-23})$ has $h=3$, with Hilbert class field generated by the famous cubic
$x^3-x-1$. The three reduced forms give $\tau_1=(1+i\sqrt{23})/2$ and
$\tau_{2,3}=(\pm1+i\sqrt{23})/4$; $G_4(\tau_1)$ is real and $G_4(\tau_{2,3})$ form a
conjugate pair.

\begin{result}[disc $-23$ class product]\label{res:d23}
\dps{canary $1.4\times10^{-160}$}
\[
G_4(\tau_1)\,G_4(\tau_2)\,G_4(\tau_3)
=\frac{11\cdot17\;P^4}{2^{24}\cdot3^6\cdot23^5\cdot\pi^{16}},\qquad
P=\!\!\prod_{\substack{k\bmod 23\\ (\frac{k}{23})=+1}}\!\!\Gamma\!\Bigl(\frac{k}{23}\Bigr)
\quad\text{(11 factors)}.
\]
\end{result}

\begin{result}[the $\gamma_2$ class polynomial from row-sums]\label{res:d23poly}
With $\gamma_2=E_4/\eta^8$ phase-corrected by the canonical cube root (the $q^{1/3}$
ambiguity of $\eta^8$ at these points), the three values $\gamma_2(\tau_i)$ computed from the
polygamma row-sums are exactly the roots of the classical class polynomial
\[
x^3+155x^2+650x+23375=0
\qquad\text{\dps{coefficients integer to $10^{-80}$}},
\]
whose constant term $23375=5^3\cdot11\cdot17$ forecast the primes $11$, $17$ in the class
product --- the norm heuristic of Section~\ref{sec:h2}, one degree up.
\end{result}

\begin{negative}[weight-4 sums are not Galois-clean]\label{neg:sums}
The dimensionless symmetric functions $e_1^3/e_3$ and $e_2^3/e_3^2$ of the raw $G_4$ values
are \emph{not} rational (PSLQ to coefficient height $10^{30}$ at $300$ dps). Unlike the value
\emph{product} (a norm) and the weight-$0$ $\gamma_2$ invariants, sums of weight-$4$ CM
values across ideal classes are not Galois-stable --- the period correction under
$\operatorname{Gal}(H/K)$ is nontrivial at weight $4$. Class polynomials must be formed from
$\gamma_2=E_4/\eta^8$ (or $j$), not from the $G_4$ values themselves.
\end{negative}

\section{Class number 4: discriminant \texorpdfstring{$-39$}{-39} and genus structure}
\label{sec:h4}

The four reduced forms of discriminant $-39$, namely $(1,1,10)$, $(2,\pm1,5)$, $(3,3,4)$,
give $\tau\in\{(1+i\sqrt{39})/2,\ (\pm1+i\sqrt{39})/4,\ (3+i\sqrt{39})/6\}$.

\begin{result}[disc $-39$ class product]\label{res:d39}
\dps{$160$-dps canary}
\[
\prod_{i=1}^{4}G_4(\tau_i)
=\frac{17\cdot23\cdot29\;P^4}{2^{28}\cdot3^{9}\cdot5^{4}\cdot13^{8}\cdot\pi^{16}},\qquad
P=\!\!\prod_{\substack{k\bmod 39\\ \chi_{-39}(k)=+1}}\!\!\Gamma\!\Bigl(\frac{k}{39}\Bigr)
\quad\text{(12 factors)}.
\]
\end{result}

A wrinkle relative to Section~\ref{sec:h3}: $\gamma_2$ is \emph{not} a class invariant here
($3\mid 39$, the classical obstruction \cite{Cox}; observed directly --- no cube-root phase choice makes
the $\gamma_2$-sums rational). But $j$ always is, and the row-sum values deliver the Hilbert
class polynomial outright.

\begin{result}[Hilbert class polynomial of disc $-39$]\label{res:H39}
\[
H_{-39}(x)=x^4+331531596\,x^3-429878960946\,x^2+109873509788637459\,x
+20919104368024767633
\]
\dps{$e_1,\dots,e_3$ PSLQ-exact integers; $e_4$ integral to $10^{-137}$}.
\end{result}

Worth dwelling on: a quartic with $20$-digit integer coefficients recovered from sums of
polygamma values at quadratic-irrational complex points. The polygamma row-sum is a
perfectly serviceable special-function front door to the classical CM tower.

\begin{result}[genus structure inside the quartet]\label{res:genus}
The class group is $\mathbb{Z}/4$, with the ambiguous form $(3,3,4)$ as the order-$2$ class.
The raw cross-class ratio $v_1/v_4$ is honestly \emph{quartic} (no degree-$2$ minimal
polynomial to height $10^{12}$), but the genus-wise ratio --- principal genus $\{1,A^2\}$
over $\{A,A^3\}$ --- is exactly quadratic, in the real part of the genus field
$\mathbb{Q}(\sqrt{-3},\sqrt{13})$ as genus theory demands
\dps{$160$-dps canary; $\sqrt3$, $\sqrt{39}$ declined by PSLQ}:
\[
\frac{G_4(\tau_1)\,G_4(\tau_4)}{G_4(\tau_2)\,G_4(\tau_3)}
=\frac{103299+4440\sqrt{13}}{181424},\qquad
N=\Bigl(\frac{9}{16}\Bigr)^{\!2},\qquad 181424=16\cdot17\cdot23\cdot29 .
\]
The denominator reuses the class-product primes $17$, $23$, $29$.
\end{result}

\section{Class number 5: discriminant \texorpdfstring{$-47$}{-47}}\label{sec:h5}

The five reduced forms $(1,1,12)$, $(2,\pm1,6)$, $(3,\pm1,4)$ give the quintet
$\tau\in\{(1+i\sqrt{47})/2,\ (\pm1+i\sqrt{47})/4,\ (\pm1+i\sqrt{47})/6\}$.

\begin{result}[disc $-47$ class product]\label{res:d47}
\dps{$200$-dps canary}
\[
\prod_{i=1}^{5}G_4(\tau_i)
=\frac{11^2\cdot23\cdot29\;P^4}{2^{52}\cdot3^{6}\cdot47^{9}\cdot\pi^{36}},\qquad
P=\!\!\prod_{\substack{k\bmod 47\\ (\frac{k}{47})=+1}}\!\!\Gamma\!\Bigl(\frac{k}{47}\Bigr)
\quad\text{(23 factors)}.
\]
\end{result}

\begin{result}[Hilbert class polynomial of disc $-47$]\label{res:H47}
The row-sum $j$-values deliver the quintic Hilbert class polynomial
\dps{all symmetric functions integral to $10^{-129}$ or better}:
\begin{align*}
H_{-47}(x)=x^5&+2257834125\,x^4-9987963828125\,x^3+5115161850595703125\,x^2\\
&-14982472850828613281250\,x+16042929600623870849609375 .
\end{align*}
\end{result}

The CM tower is now polygamma-explicit through class number $5$ with one uniform recipe:
\emph{row-sum $\to$ class-product log-PSLQ (norm-forecast primes) $\to$
$j$-symmetric-function integrality.}

\section{The opaque components}\label{sec:opaque}

The components \emph{not} fixed by reflection-conjugation or by the even-row mechanism ---
$\operatorname{Im}\psi^{(1)}(i)$, $\operatorname{Re}\psi(i)$ beyond $\gamma$,
$\operatorname{Im}\psi^{(1)}\bigl(\tfrac{1+i}{2}\bigr)$ --- are \emph{odd-row} lattice data:
for example
\[
\operatorname{Im}\psi^{(1)}(i)=-2\sum_{k\ge1}\frac{k}{(k^2+1)^2}
\]
is an odd-power analogue of the $\coth$ family, which the bilateral (even) cotangent
expansion cannot reach. These are Eichler-integral/quasi-modular territory --- the
$E_2$-flavoured layer; compare $\sum_{k\ge1}k/(e^{2\pi k}-1)=\tfrac1{24}-\tfrac1{8\pi}$,
which is $E_2(i)=3/\pi$ in Lambert-series clothing.

\begin{negative}\label{neg:opaque}
A PSLQ basket for these constants against products drawn from
\[
\{1,\ \pi,\ \coth\pi,\ \pi^2/\sinh^2\pi,\ \pi\coth\pi,\ \dots\}
\]
found nothing at small coefficient heights --- consistent with genuinely new constants.
Recorded as an honest negative; the Lambert-series face ($\sum\sigma_5(n)q^n$ at
$q=e^{-2\pi}$ and $q=-e^{-\pi\sqrt3}$) is the natural next probe surface.
\end{negative}

\section{Methodology notes}\label{sec:method}

Two failure chains from the discovery process are worth recording.

\paragraph{A vanishing identity passing is weak evidence for the machinery around it.}
The row-sum prefactor in \eqref{eq:rowsum} is $2/(2k-1)!$; an early weight-$6$ probe used
$1/120$ instead of $1/60$. The \emph{vanishing} tests are blind to such a factor (zero
swallows it) --- they passed while the value probe quietly computed garbage. The bug was
caught by cross-checking the lattice row-sum against the modular $q$-series
$G_6=2\zeta(6)E_6(\rho)$ with $q=-e^{-\pi\sqrt3}$ (agreement to $30$+ digits).

\paragraph{When a CM-value PSLQ fails, anchor on $\eta$.}
With the correct value in hand, log-PSLQ over $\{\ln\Gamma(1/3),\ln\pi,\ln2,\ln3\}$ still
found nothing for $G_6(\rho)$ --- a span failure, not a precision failure:
$8960=2^8\cdot5\cdot7$ and $\zeta(6)=\pi^6/945$ drag in the primes $5$ and $7$. The
debugging anchor was the Dedekind eta function: PSLQ over the same small basis instantly
returned the classical $|\eta(\rho)|=3^{1/8}\Gamma(1/3)^{3/2}/(2\pi)$, and then
$E_6(\rho)^2=1728\,|\eta(\rho)|^{24}$ forces the closed form with no further search
(note $\eta(\rho)^{24}=\Delta(\rho)$ is \emph{negative} real here, since
$q=-e^{-\pi\sqrt3}$ and $E_4(\rho)=0$). Eta's
closed forms have tiny prime support; derive upward through the modular algebra rather than
widening a PSLQ basis blindly.

\section{Open directions}\label{sec:open}

\begin{enumerate}
\item The $G_6$ row at the class-number $\ge2$ points, and the finer Galois structure inside
the disc $-39$ quartet (which pairs of $G_4$ values generate which subfield of the class
field).
\item The quasi-modular layer: $E_2(i)=3/\pi$ row-summed is weight $2$ --- conditionally
convergent rows requiring Eisenstein summation --- and should explain exactly which
combinations of the opaque odd constants of Section~\ref{sec:opaque} are reachable.
\item Folding the period identities into the repository's identity stores once a store schema
for infinite-sum claims exists (current stores hold finite reduction rules only), and
converting the finite-rank discoveries into multiplier proof certificates.
\end{enumerate}

\section*{Sources and verification artifacts}

This article presents results recorded in the Smithereens repository
(\texttt{src/PolyLog}), principally in the report
\path{docs/reports/polygamma-complex-arguments-cm-lattices__d41f7be20c61.md} and the
commit range \texttt{272af472..68236dab4} (2026-06-09 through 2026-06-11). All numerical
verification used \texttt{mpmath} at the per-identity precisions quoted in square brackets;
PSLQ discoveries were re-verified at elevated precision (canary discipline). The classical
ingredients --- the Chowla--Selberg formula, the structure of $M_*(\mathrm{SL}_2(\mathbb{Z}))$,
Hurwitz's recursion, genus theory, and the class polynomials of discriminants $-23$, $-39$,
$-47$ --- are classical; the polygamma-explicit formulations and the specific
ratio/class-product evaluations are, to the project's knowledge, newly recorded here.

\begin{thebibliography}{9}

\bibitem{ChowlaSelberg}
A.~Selberg and S.~Chowla, \emph{On Epstein's zeta-function}, J.~reine angew.\ Math.\
\textbf{227} (1967), 86--110.

\bibitem{Damerell}
R.~M.~Damerell, \emph{$L$-functions of elliptic curves with complex multiplication, I},
Acta Arith.\ \textbf{17} (1970), 287--301.

\bibitem{Hurwitz}
A.~Hurwitz, \emph{\"Uber die Entwicklungskoeffizienten der lemniskatischen Funktionen},
Math.\ Ann.\ \textbf{51} (1899), 196--226.

\bibitem{Cox}
D.~A.~Cox, \emph{Primes of the Form $x^2+ny^2$: Fermat, Class Field Theory, and Complex
Multiplication}, Wiley, 1989. (Genus theory; class invariants; the $3\mid D$ obstruction
for $\gamma_2$.)

\bibitem{MathWorldDigamma}
E.~W.~Weisstein, \emph{Digamma Function}, MathWorld --- A Wolfram Web Resource.
(Source of the FoxTrot identity.)

\end{thebibliography}

\end{document}
